\chapter{Metrik Tamlık ve Kompaktlık}

Metrik tamlık (completeness), bir metrik uzayda Cauchy dizilerinin yakınsayıp
yakınsamadığını araştıran temel bir kavramdır.

%-------------------------------------------------------------------
\section{Konu}

\subsection{Cauchy Dizisi ve Tamlık}

\begin{definition}[Cauchy Dizisi]
$(M, d)$ metrik uzayında $\{x_n\}$ dizisi \textbf{Cauchy} ise:
\[
\forall\varepsilon>0,\;\exists N \in \mathbb{N}:\;
m,n \geq N \Rightarrow d(x_m, x_n) < \varepsilon.
\]
\end{definition}

\begin{definition}[Tam Metrik Uzay]
$(M, d)$ \textbf{tam} ise: Her Cauchy dizisi $M$'de bir limite yakınsır.
\end{definition}

\subsection{Tamamen Sınırlılık}

$A \subseteq M$ \textbf{tamamen sınırlı} ise: her $\varepsilon>0$ için $A$'nın
sonlu bir $\varepsilon$-net kümesi vardır.
\[
S = \{s_1,\ldots,s_n\} \subseteq M,\quad
\forall x \in A\;\exists i:\; d(x, s_i) < \varepsilon.
\]

\subsection{Metrik Kompaktlık Karakterizasyonu}

\[
(M, d) \text{ tam} \wedge \text{tamamen sınırlı}
\Longleftrightarrow
(M, d) \text{ kompakt.}
\]

\begin{table}[h]
\centering
\begin{tabular}{llll}
\toprule
\textbf{Uzay} & \textbf{Tam?} & \textbf{Tam. sınırlı?} & \textbf{Kompakt?} \\
\midrule
$\mathbb{R}$ & \checkmark & $\times$ & $\times$ \\
$[0,1]$ & \checkmark & \checkmark & \checkmark \\
$(0,1)$ & $\times$ & \checkmark & $\times$ \\
$\mathbb{Q}$ & $\times$ & $\times$ & $\times$ \\
Sonlu metrik & \checkmark & \checkmark & \checkmark \\
\bottomrule
\end{tabular}
\caption{Tamlık ve kompaktlık örnekleri}
\end{table}

%-------------------------------------------------------------------
\section{Teoremler}

\begin{theorem}[Heine-Borel Metrik Genelleştirmesi]
Metrik uzayda kompaktlık $\Leftrightarrow$ tamlık $\wedge$ tamamen sınırlılık.
\end{theorem}

\begin{theorem}[Banach Sabit-Nokta Teoremi]
$(M, d)$ tam metrik uzay; $T: M \to M$ büzülme ($K < 1$)
$\Rightarrow$ $T$'nin eşsiz sabit noktası $x^*$ vardır: $T(x^*) = x^*$.
\end{theorem}

\begin{theorem}[Baire Kategorisi Teoremi]
Tam metrik uzay 1.\ kategoriden değildir.
\end{theorem}

%-------------------------------------------------------------------
\section{Algoritmalar}

\subsection{Sonlu Metrikte is\_complete ve is\_totally\_bounded}

Sonlu metrik uzayda her Cauchy dizisi eninde sonunda sabit: trivially tam.
Taşıyıcı kendisi $\varepsilon$-net: trivially tamamen sınırlı.
Her ikisi: $O(1)$.

\subsection{metric\_compactness\_check}

\begin{algorithm}
\caption{MetricCompactnessCheck$(M, d)$}
\begin{algorithmic}[1]
\State complete $\leftarrow$ is\_complete$(M, d)$
\State totally\_bounded $\leftarrow$ is\_totally\_bounded$(M, d)$
\State \Return complete $\wedge$ totally\_bounded
\end{algorithmic}
\end{algorithm}

%-------------------------------------------------------------------
\section{pytop API}

\begin{lstlisting}[language=Python]
from pytop.metric_completeness import (
    is_complete,
    is_totally_bounded,
    metric_compactness_check,
    analyze_metric_completeness,
)
\end{lstlisting}

%-------------------------------------------------------------------
\section{Örnekler}

\subsection*{Örnek 5.1 — Sonlu Metrik: Tam ve Kompakt}

\begin{lstlisting}[language=Python]
points = ['A', 'B', 'C', 'D']
dist = {(a, b): (0 if a == b else 1)
        for a in points for b in points}
fms = FiniteMetricSpace(carrier=tuple(points), distance=dist)
print("is_complete:", is_complete(fms).status)
print("is_totally_bounded:", is_totally_bounded(fms).status)
mc = metric_compactness_check(fms)
print("metric_compactness:", mc.status, mc.value)
\end{lstlisting}

\begin{verbatim}
is_complete: true
is_totally_bounded: true
metric_compactness: true True
\end{verbatim}

\subsection*{Örnek 5.4 — analyze\_metric\_completeness}

\begin{lstlisting}[language=Python]
r = analyze_metric_completeness(fms)
print("status:", r.status)
print("value:", r.value)
\end{lstlisting}

\begin{verbatim}
status: true
value: {'is_complete': True, 'is_totally_bounded': True, 'metric_compact': True}
\end{verbatim}

%-------------------------------------------------------------------
\section{Alıştırmalar}

\subsection*{Kodlama}

\begin{enumerate}[label=K\arabic*.]
\item Rasyonel sayılar $\mathbb{Q}$ (sembolik) için \texttt{is\_complete} sonucunu
      kontrol edin.
\item Kendi 4-noktalı metrik uzayınızı oluşturun ve \texttt{metric\_compactness\_check}
      uygulayın.
\item \texttt{analyze\_metric\_completeness} çalıştırın ve \texttt{value} sözlüğünü inceleyin.
\end{enumerate}

\subsection*{Teori}

\begin{enumerate}[label=T\arabic*.]
\item Banach sabit-nokta teoremini sözlü olarak açıklayın ve bir uygulama verin.
\item Kapalı $[0,1]$'in tam, açık $(0,1)$'in tam olmadığını gösterin.
      (Hint: $1/n$ dizisi Cauchy ama $0 \notin (0,1)$.)
\end{enumerate}
